%------------------------------------------------------------------------------
% 3.7 ملخص الفصل — rewritten (2026-10) to reflect:
%   sprints with Analyst/Reporter in S3 + Onboarding/Audit in S4,
%   29 functional requirements, 7 use cases, two-level human actor taxonomy,
%   and the NEW agent-modeling section (3.5).
%------------------------------------------------------------------------------
\section{ملخص الفصل}
\label{sec:analysis-summary}

{\large
استعرض هذا الفصل البنية المنهجية والتحليلية الشاملة التي استند إليها مشروع المساعد المؤسسي الذكي \textbf{«بيان»}، مشكلاً حجر الأساس قبل الانتقال إلى مرحلة التصميم الفعلي. وقد تم تناول المحاور الرئيسية التالية:

\begin{itemize}
    \item \textbf{منهجية التطوير:} توضيح دورة التطوير الرشيقة (\EN{Agile}) المعتمدة لبناء النظام، والتي قُسّمت إلى أربع دورات تطويرية تضمن المرونة والتسليم التدريجي: بناء وكيل \EN{SQL} وطبقة الرؤى الذكية، فوكيل \EN{RAG} والفهرسة، ثم اكتمال منظومة الوكلاء الخمسة ببناء وكيلي التحليل والتقارير وتنسيقها عبر الوكيل المشرف ودمجها بالواجهة العربية، وأخيراً التهيئة الذكية وسجل التدقيق والاختبار الشامل والنشر.

    \item \textbf{دراسة الجدوى:} تقديم دراسة مستفيضة شملت الأبعاد الاقتصادية، والتقنية، والتشغيلية، والقانونية، والمؤسسية، مما أثبت قابلية المشروع للتطبيق العملي ونجاعته في تلبية احتياجات المؤسسات المستهدفة.

    \item \textbf{تحليل المتطلبات:} تحديد تسعة وعشرين متطلباً وظيفياً موزعة على خمس فئات وفق أولويات \EN{MoSCoW} (20 عالية، و8 متوسطة، ومتطلب منخفض واحد)، وعشرين متطلباً غير وظيفي وفق معيار \EN{ISO/IEC 25010}، إضافة إلى متطلبات بيئة التشغيل وقيود النظام وافتراضاته.

    \item \textbf{حالات الاستخدام (\EN{Use Cases}):} نمذجة تفاعلات المستخدمين مع النظام عبر سبع حالات استخدام (\EN{UC-01} إلى \EN{UC-07})، وفق معايير \EN{UML} التي تقصر تمثيل الفاعلين على البشريين خارج حدود النظام: مسؤول المنصة (\EN{PA})، ومستخدمو المؤسسة (\EN{TU}) المتفرعين بالتعميم إلى مسؤول تقنية المعلومات والمدير والموظف، مع نمذجة الوكلاء مكوناتٍ داخلية للنظام.

    \item \textbf{نمذجة منظومة الوكلاء والعلاقات التفاعلية بينهم:} توصيف الوكلاء الخمسة هيكلياً (الدور والمدخلات والأدوات والمخرجات)، وتحديد بروتوكول التخاطب بينهم عبر كائن الحالة المشترك \EN{AgentState} وحقوله الاثني عشر، وعرض أنماط التنسيق الأربعة (التفويض الأحادي، والسلسلة التحليلية المتتابعة، والتكامل الهجين بين اللوائح والبيانات، وحلقة التصحيح الذاتي)، وتوثيق التسلسل التفاعلي الشامل بمخطط تتابع \EN{UML} يحقق المهمة المجدولة في الخطة الزمنية للفصل الأول.

    \item \textbf{تحليل المخاطر:} استعراض سجل المخاطر الشامل ومصفوفة تقييمها، مع وضع خطط استجابة محكمة للتعامل مع المخاطر عالية الأولوية والحرجة لضمان استمرارية سير العمل بنجاح.
\end{itemize}

يوفر هذا الفصل الرؤية التحليلية المتكاملة والمتطلبات الدقيقة وبروتوكول التفاعل الداخلي للمنظومة، ممهداً الطريق للانتقال إلى الفصل الرابع المعني بتصميم المعمارية الهيكلية، ومنظومة الوكلاء وطبقة الرؤى الذكية، وقواعد البيانات، وإطار الأمان والصلاحيات، وواجهات النظام.
}
%------------------------------------------------------------------------------
